% !TeX root = ../main.tex
\section{Recovery Costs: Recomputation or Restoration}
\label{sec:tradeoff}

Better scheduling reduces cache churn but does not eliminate recovery after eviction.
We extend the replication by studying recovery when useful state no longer fits on the GPU: discard it for later recomputation, or preserve it in host memory for restoration.
For $n$ missing prefix tokens with KV size $S$, a serial save-and-restore cycle costs
\begin{equation}
 T_{\mathrm{off}}(S)=T_{\mathrm{out}}(S)+T_{\mathrm{in}}(S)+T_{\mathrm{mgr}}.
 \label{eq:offload-cost}
\end{equation}
The terms cover preservation, restoration, and management such as allocation and lookup.
Offloading is cheaper in elapsed service time when this sum is below conditional recomputation time $T_{\mathrm{rec}}(n)$.
With one possible reuse of probability $p$, the expected-work condition becomes
\begin{equation}
 T_{\mathrm{out}}+pT_{\mathrm{in}}+T_{\mathrm{mgr}}<pT_{\mathrm{rec}}.
 \label{eq:expected-offload-cost}
\end{equation}
This model assumes one write, at most one reuse, and no overlap or contention.
Clean retained copies can amortize writes over further reuse; speculative saves that are never reused can waste work.

The implemented host tier pays device-to-host and host-to-device copy costs.
A future SSD path must additionally pay storage writes, reads, staging, and queueing; host bandwidth alone cannot establish its break-even point.
The following study reports end-to-end improvements at selected host-DRAM operating points, explicitly adding backing capacity.
It does not isolate prefill time, establish equal-total-memory gains, or measure energy and monetary cost.
